\section{Exact centering and the hard-sphere fluctuation limit}
\label{sec:transfer}

The pasted Gaussian limit and the dynamical comparison have now been
proved independently.  Choose \(\Gamma\) large enough for
\Cref{thm:marked-expansion,prop:cutoff}, and then fix \(L\) large
enough for those results and \Cref{thm:sharp-cumulants}.  Their
thresholds depend only on the fixed kinetic data and \(\Gamma\),
not on the observation list or moment order.  Thus all three results
apply to one pasted dynamics.  The remaining step is to transfer its
exact centering, since equality outside an event of small probability
alone would not compare expectations at the fluctuation scale.

\begin{lemma}[Exact-centering comparison]\label{lem:center-transfer}
Let the true and pasted dynamics start from the same initial ensemble.
Under the cutoff estimate of \Cref{prop:cutoff}, the expectation of
any fixed bounded empirical observation list differs between the two
dynamics by \(o(\mu^{-1/2})\).  Consequently the corresponding
exact-centered fluctuation vectors have the same weak limits.
\end{lemma}

\begin{proof}
Let \(\mathcal E_\eps\) be the event that the two trajectories separate
at some time in \([0,T]\).  By \Cref{prop:cutoff},
\(\P(\mathcal E_\eps)=O(\eps^{1+\eta})\) for some \(\eta>0\).
Both dynamics preserve \(N\), and their observations coincide on its
complement.  For the scalar record sum, denoted by \(S^{\eps,\mathrm{tr}}\)
and \(S^\eps\) on the two dynamics,
\[
 |S^{\eps,\mathrm{tr}}-S^\eps|
       \le 2M_{\mathrm{obs}}N\ind_{\mathcal E_\eps}.
\]
Choose an integer \(p>(1+\eta)/\eta\).  By Hölder's inequality and
\eqref{eq:initial-N-moments},
\begin{equation}\label{eq:limit-center-estimate}
 \begin{split}
 \frac{|\E S^{\eps,\mathrm{tr}}-\E S^\eps|}{\sqrt\mu}
 &\le 2M_{\mathrm{obs}}\sqrt\mu
       \bigl(\E(N/\mu)^p\bigr)^{1/p}
            \P(\mathcal E_\eps)^{1-1/p}\\
 &\le C_p\,\eps^{-1+(1+\eta)(1-1/p)}
       \longrightarrow0.
 \end{split}
\end{equation}
The uncentered sums divided by \(\sqrt\mu\) differ only on an event
whose probability tends to zero, so their difference tends to zero
in probability.  Equation \eqref{eq:limit-center-estimate} gives the
same conclusion after subtracting their respective exact
expectations.  This argument applies to each component of a finite
list, or directly to all its linear combinations.  It requires no
bound on high moments of the scaled fluctuation restricted to
\(\mathcal E_\eps\).
\end{proof}

\begin{proof}[Proof of \Cref{thm:main}]
The Gaussian generalized field, its initial covariance, its noise
normalization, and its well-posedness have been established in
\Cref{prop:gaussian-target}.  Fix an arbitrary real linear combination
of the finitely many empirical observations in the theorem, and use
its coefficients in the mark updates \eqref{eq:limit-mark-update}.
The exact identity \eqref{eq:limit-record-identity},
\Cref{prop:cumulant-clt}, and \Cref{lem:center-transfer} give convergence
of this linear combination, with the exact microscopic centering, to
the corresponding linear combination of the Gaussian observations.
Its variance is \eqref{eq:limit-variance-formula}, hence is the variance
specified by \eqref{eq:target-covariance}.  Since this holds for every
real choice of coefficients, the Cramér--Wold theorem proves
\eqref{eq:main-convergence}.
\end{proof}
